\chapter{Beta-2 Microglobulin}

\section*{What Is This Test?}
Beta-2 microglobulin (B2M) is an 11.8-kDa protein. It is the invariable beta-chain (light chain) of the major histocompatibility complex class I (MHC class I) molecule. This molecule is present on the surface of all nucleated cells. B2M binds non-covalently to the alpha (heavy) chain of the MHC class I molecule (HLA-A, HLA-B, and HLA-C). B2M is essential for stable cell-surface expression of MHC class I. This molecule presents intracellular peptide antigens to CD8+ cytotoxic T lymphocytes. B2M is shed continuously from the cell surface during normal MHC class I turnover. It is freely filtered through the glomerular basement membrane. Under normal conditions, >99\% of filtered B2M is reabsorbed and catabolized in the renal proximal tubules. Serum B2M therefore reflects the balance between production and renal elimination. Production comes primarily from lymphoid cells. Serum B2M is high in conditions of increased lymphocyte turnover. These include hematologic malignancies, chronic inflammation, and autoimmune disease. Serum B2M is also high in decreased renal clearance. These include renal tubular dysfunction and chronic kidney disease. B2M is a key prognostic marker and staging criterion in multiple myeloma. It is used in the International Staging System (ISS). It is also a prognostic marker in chronic lymphocytic leukemia (CLL), diffuse large B-cell lymphoma (DLBCL), Hodgkin lymphoma, and myelodysplastic syndromes.

\section*{Alternative Names}
\begin{itemize}
  \item B2M
  \item Beta-2-Microglobulin
  \item $\beta$2-Microglobulin
  \item $\beta$2M
  \item Serum Beta-2 Microglobulin
  \item Beta-2-M
  \item Thymotaxin
\end{itemize}
\section*{Principle}
The laboratory measures B2M by particle-enhanced immunoturbidimetry or immunonephelometry. These methods run on automated clinical chemistry analyzers. B2M can also be measured by chemiluminescent immunoassay (CLIA) on dedicated immunoassay platforms. Latex particles coated with anti-human B2M antibodies are incubated with patient serum or plasma. B2M in the sample binds to the antibody-coated particles. This forms immune complexes that cause agglutination. The resulting increase in turbidity (light scattering) is measured photometrically. It is proportional to the B2M concentration. Results are calibrated against an international reference standard (WHO 1st IS for B2M, 85/507). They are reported in mg/L (or mcg/mL, which is numerically equivalent). B2M is unstable in acid conditions. It degrades at room temperature over extended time periods. Serum specimens should be processed promptly. Urine B2M can be measured at the same time as serum to evaluate renal tubular function. B2M is freely filtered, and >99\% is reabsorbed in the proximal tubules. Therefore, increased urinary B2M excretion with normal serum B2M and normal GFR indicates a proximal tubular defect. Examples are Fanconi syndrome, heavy metal toxicity, and drug-induced tubulopathy.

\section*{History}
B2M was first isolated in 1964 by Berggard and Bearn \cite{berggard1968isolation}. They isolated it from the urine of patients with Wilson disease and chronic cadmium poisoning (renal tubular proteinuria). Its amino acid sequence was determined in 1972. Shortly after, its structural homology with immunoglobulin domains was recognized. Specifically, it matched the CH3 domain of IgG. This led to the identification of B2M as part of the MHC class I molecule. The role of B2M in multiple myeloma prognostication was recognized in the 1980s. It was incorporated into the Durie-Salmon staging system (1975). Later it was added to the International Staging System (ISS) in 2005 \cite{greipp2005international}. The link between B2M and dialysis-related amyloidosis (DRA) was established in the 1980s. In end-stage renal disease, B2M accumulates due to impaired renal clearance. B2M amyloid fibrils deposit in joints, bones, and the carpal tunnel.

\section*{Why This Test Is Ordered}
\begin{itemize}
  \item Staging and prognostication of multiple myeloma --- B2M is a key component of the Revised International Staging System (R-ISS) for multiple myeloma. B2M <3.5 mg/L and albumin $\geq$3.5 g/dL = stage I (median survival 62 months). B2M $\geq$5.5 mg/L = stage III (median survival 29 months). The R-ISS (2015) additionally incorporates serum LDH and high-risk cytogenetics [del(17p), t(4;14), t(14;16)] \cite{palumbo2015revised}.
  \item Prognostic assessment in chronic lymphocytic leukemia (CLL) --- high B2M is an independent adverse prognostic factor in CLL. It is linked to shorter time to first treatment and shorter overall survival. B2M is incorporated into prognostic scores including the CLL-International Prognostic Index (CLL-IPI).
  \item Prognostication in diffuse large B-cell lymphoma (DLBCL) --- high B2M at diagnosis is an independent adverse prognostic factor. It is integrated into some prognostic scores (SELL, NCCN-IPI).
  \item Monitoring of disease activity and response to therapy in multiple myeloma, where B2M levels reflect tumor burden. Falling B2M indicates treatment response. Rising B2M indicates disease progression.
  \item Evaluation of renal tubular function --- simultaneous measurement of serum B2M and urine B2M can distinguish between glomerular disease and proximal tubular dysfunction. Glomerular disease shows high serum B2M with modest urine B2M. Proximal tubular dysfunction shows normal or slightly high serum B2M with markedly high urine B2M. Urine B2M >0.3 mg/L suggests proximal tubular injury.
  \item Prognostic assessment in myelodysplastic syndromes (MDS), Hodgkin lymphoma, and aggressive non-Hodgkin lymphomas.
  \item Monitoring for dialysis-related amyloidosis (DRA) in long-term hemodialysis patients --- persistently high B2M >20 mg/L in end-stage renal disease is a risk factor for B2M amyloid deposition in joints and carpal tunnel.
\end{itemize}

\section*{Primary vs Secondary Test}
B2M is a secondary prognostic test used with other clinical, laboratory, and cytogenetic factors in multiple myeloma and lymphoproliferative disorders. It is a required component of the International Staging System for multiple myeloma.

\section*{How This Test Is Performed}
For most patients, a health care professional takes a blood sample from a vein in your arm using a small needle. After the needle is inserted, a small amount of blood is collected into a test tube or vial. You may feel a little sting when the needle goes in or out. This usually takes less than five minutes. For concomitant urine B2M measurement, a random or 24-hour urine collection may also be required.

\section*{What Preparation Is Needed}
No special preparation is required. Fasting is not necessary. For urine B2M measurement, avoid vigorous exercise for 24 hours before collection, as heavy exercise can cause transient tubular proteinuria. B2M degrades in acidic urine (pH <5.5). The patient may be instructed to take sodium bicarbonate to alkalinize the urine before collection. This prevents false-negative urinary B2M results.

\section*{Contraindications and Precautions}
\begin{itemize}
  \item No absolute contraindications. The most critical limitation of serum B2M is its dependence on renal function. Any condition that impairs glomerular filtration rate (GFR) raises serum B2M, regardless of the underlying disease. Serum B2M must ALWAYS be interpreted with the patient's serum creatinine, eGFR, and clinical status. In end-stage renal disease (GFR <15 mL/min), B2M may reach 30--40 mg/L due to impaired renal clearance alone, without any hematologic malignancy. In acute kidney injury (AKI) or chronic kidney disease (CKD), B2M is raised in proportion to the degree of renal impairment. Conditions other than hematologic malignancies can also raise B2M. These include autoimmune diseases (systemic lupus erythematosus, rheumatoid arthritis, Sjogren syndrome, sarcoidosis). They include chronic viral infections (HIV/AIDS, chronic hepatitis B and C, CMV). They include acute organ transplant rejection, lymphoproliferative disorders other than myeloma (Hodgkin lymphoma, non-Hodgkin lymphomas), and reactive lymphoid hyperplasia. Collect the specimen in a serum separator tube (SST, gold-top) or red-top tube. B2M degrades at room temperature. Serum should be separated within 4 hours of collection. Hemolysis and lipemia can interfere with immunoturbidimetric assays.
\end{itemize}
\section*{Risks}
Minimal risk. Slight pain or bruising at the venipuncture site. Rare: hematoma, infection, vasovagal syncope.

\section*{What Happens After the Test}
Results are typically available within 1 to 2 hours. B2M results must be interpreted with serum creatinine/eGFR, serum albumin, serum protein electrophoresis (SPEP), immunofixation, serum free light chains, and the clinical picture. In the workup of multiple myeloma, B2M is required for ISS staging. It should be obtained at diagnosis. Serial B2M measurements may be used to monitor treatment response and disease progression. Serial serum free light chain and M-protein measurements are the primary monitoring tools.

\section*{Understanding Results}

\subsection*{Below Normal}
\begin{itemize}
  \item B2M <1.5 mg/L (or <1.0 mg/L, depending on age and laboratory): Normal. Indicates low tumor burden or inactive disease. In multiple myeloma, low B2M with normal renal function is a favorable prognostic factor (ISS Stage I with albumin $\geq$3.5 g/dL).
  \item Decreasing B2M after initiation of chemotherapy or novel agent therapy for multiple myeloma: Indicates treatment response. A reduction of $\geq$50\% in B2M from the pretreatment value is a favorable sign. It correlates with disease response and prolonged progression-free survival.
  \item Normal B2M after allogeneic stem cell transplantation for myeloma or lymphoma: Indicates disease control and favorable long-term prognosis.
  \item Normal B2M in the setting of lymphadenopathy or lymphocytosis: Argues against a high-grade lymphoproliferative disorder. It does not exclude indolent lymphoma or early-stage CLL.
\end{itemize}

\subsection*{Above Normal}
\begin{itemize}
  \item B2M 2.0--3.4 mg/L (mild elevation): Mild elevation may reflect early renal impairment or low-grade lymphoid proliferation. In multiple myeloma, this range with albumin $\geq$3.5 g/dL is consistent with ISS Stage I. In CLL, mild elevation is a modest adverse prognostic factor.
  \item B2M 3.5--5.5 mg/L (moderate elevation): This range in multiple myeloma defines ISS Stage II (intermediate prognosis) when renal function is normal. In CLL and DLBCL, moderate B2M elevation is an adverse prognostic factor. It is linked to shorter overall survival.
  \item B2M >5.5 mg/L (marked elevation): In multiple myeloma, this range defines ISS Stage III (poor prognosis). In the ISS 2005, median survival is 29 months. In the R-ISS, Stage III with normal LDH and standard-risk cytogenetics has a median OS of 43 months. B2M >5.5 mg/L with high LDH and high-risk cytogenetics defines R-ISS Stage III (median OS 43 months). This degree of elevation correlates with high tumor burden. It should prompt urgent initiation of therapy.
  \item B2M >10 mg/L: Severe elevation. In multiple myeloma, this indicates a very high tumor burden and/or significant renal impairment from myeloma cast nephropathy. In chronic lymphocytic leukemia, B2M >4 mg/L is linked to advanced-stage disease (Rai III--IV, Binet C) and short overall survival.
  \item B2M elevation from renal impairment: Any cause of renal dysfunction (AKI, CKD, ESRD) raises serum B2M. The B2M/creatinine ratio or correlation of B2M with eGFR can help distinguish B2M elevation due to tumor burden versus renal impairment. In ESRD on hemodialysis, B2M 20--40 mg/L is typical. It reflects impaired renal clearance alone.
  \item B2M elevation from non-malignant conditions: Autoimmune diseases (active SLE, Sjogren syndrome, rheumatoid arthritis), sarcoidosis, chronic infections (HIV/AIDS, chronic HBV/HCV), acute organ transplant rejection (B2M rises with immune activation against the allograft), and reactive lymphoid hyperplasia.
  \item B2M >15--20 mg/L in ESRD on long-term hemodialysis: Risk factor for dialysis-related amyloidosis (DRA) with B2M amyloid fibril deposition in the carpal tunnel (carpal tunnel syndrome), joints (destructive arthropathy), and bone (cystic lesions). High-flux hemodialysis membranes and hemodiafiltration can lower B2M levels.
\end{itemize}

\section*{Normal Results}
\begin{itemize}
  \item \textbf{Serum Beta-2 Microglobulin (adults, age <60 years):} 0.8--2.2 mg/L (laboratory-dependent; some use <1.5 mg/L, others <2.0 mg/L)
  \item \textbf{Serum Beta-2 Microglobulin (adults, age >60 years):} Slightly higher reference range (0.9--2.5 mg/L) due to age-related decline in GFR
  \item \textbf{Urine Beta-2 Microglobulin:} <0.3 mg/L (random) or <0.2 mg/24 hours (timed collection). Higher levels indicate proximal tubular dysfunction
  \item \textbf{Multiple Myeloma International Staging System (ISS) --- Revised ISS (R-ISS) 2015:}
  \begin{itemize}
    \item Stage I: B2M <3.5 mg/L AND albumin $\geq$3.5 g/dL AND normal LDH AND standard-risk cytogenetics (no high-risk abnormalities) --- 5-year OS 82\%
    \item Stage II: Not meeting criteria for Stage I or III --- 5-year OS 62\%
    \item Stage III: B2M $\geq$5.5 mg/L AND either elevated LDH OR high-risk cytogenetics [del(17p), t(4;14), t(14;16)] --- 5-year OS 40\%
  \end{itemize}
  \item \textbf{B2M must always be interpreted with simultaneous assessment of renal function (serum creatinine, eGFR).} B2M elevation with normal renal function is clinically significant. B2M elevation in renal failure must be interpreted with caution.
\end{itemize}

\section*{Follow-up Tests}
\begin{itemize}
  \item \textbf{Serum protein electrophoresis (SPEP), serum immunofixation (IFE), and serum free light chain (FLC) assay:} Core tests for the diagnosis and monitoring of multiple myeloma and monoclonal gammopathies. B2M is used for staging. The M-protein and free light chains are used for diagnosis, classification, and monitoring.
  \item \textbf{Serum albumin, LDH, and complete blood count:} Additional components of the R-ISS staging system for multiple myeloma and for prognostic assessment in lymphoma.
  \item \textbf{Bone marrow biopsy and aspiration:} Essential for diagnosis of multiple myeloma. Assesses plasma cell percentage (clonal plasma cells $\geq$10\% for symptomatic myeloma), morphology, and clonality by flow cytometry or IHC. FISH tests for high-risk cytogenetic abnormalities: del(17p), t(4;14), t(14;16), t(14;20), and 1q amplification (gain/amplification).
  \item \textbf{Skeletal survey (whole-body low-dose CT, PET-CT, or whole-body MRI):} To detect lytic bone lesions, osteopenia, and extramedullary plasmacytomas.
  \item \textbf{Serum creatinine and eGFR:} Essential for interpreting B2M and for identifying renal impairment from myeloma cast nephropathy.
  \item \textbf{Flow cytometry of peripheral blood and bone marrow:} For diagnosis and minimal residual disease (MRD) monitoring in CLL, lymphoma, and myeloma.
  \item \textbf{Urine protein electrophoresis (UPEP) and urine immunofixation:} To detect Bence Jones protein (monoclonal free light chains) and quantify monoclonal proteinuria in multiple myeloma and AL amyloidosis.
  \item \textbf{Genetic testing (CLL FISH panel):} Del(17p), del(11q), del(13q), trisomy 12, and analysis of immunoglobulin heavy-chain variable region (IGHV) mutational status for prognostic assessment in CLL (B2M is combined with these for the CLL-IPI score).
\end{itemize}

\section*{References}
\bibliographystyle{unsrtnat}
\bibliography{references}

\clearpage
\section*{Regulatory Status and Authorization Date}
This chapter describes a test category, not a specific commercial product. FDA, CDSCO, EU IVDR, and MHRA approval or registration dates therefore cannot be assigned without the assay manufacturer, product name, instrument, and intended use. For laboratory-developed tests, an individual agency approval date may not exist. See the Preface for the interpretation of regulatory status.

\clearpage
